%% ============================================================
%% CHIRAL CENTER THEOREM
%% Typed operadic model, comparison criteria, and scalar firewalls
%% ============================================================

\section{The chiral endomorphism operad and the chiral center theorem}
\label{sec:chiral-center-theorem}

\index{chiral center theorem|textbf}
\index{brace algebra!chiral|textbf}
\index{derived center!typed model}

The ordinary chiral centre records degree-zero endomorphisms.  Product
deformations and their obstructions live in the derived endomorphism
object
\begin{equation}\label{eq:cct-derived-centre}
 Z^{\mathrm{der}}_{\mathrm{ch}}(A)
 :=
 R\operatorname{End}_{A^e}(A)
\end{equation}
in the derived category of complete chiral \(A\)-bimodules.  This
definition fixes the object before a bar resolution, a formal-disk
model, or a geometric configuration-space model is chosen.

The first structural question is therefore precise: which explicit
cochain model computes \eqref{eq:cct-derived-centre}, and which
operations survive the comparison?  The answer has two parts.  An
actual dg operad of chiral operations supplies braces formally.
A geometric or formal-Laurent realization computes the derived centre
once a named quasi-isomorphism has been constructed.

\begin{definition}[Operadic chart package;
\ClaimStatusDefinitional]
\label{def:cct-operadic-chart-package}
For a complete augmented chiral algebra \(A\), the package
\(H_{\mathrm{op}}(A)\) consists of:
\begin{enumerate}[label=\textup{(\roman*)}]
\item a complete one-coloured dg nonsymmetric operad
 \(\mathcal P_A^{\mathrm{ch}}\) whose arity-\(n\) term models continuous
 \(n\)-ary chiral operations on \(A\), including arity zero;
\item continuous partial compositions in a fixed common expansion
 domain, with the operadic unit and associativity identities;
\item a suspension convention for which
 \[
 \mathfrak g_A
 :=
 C^\bullet_{\mathcal P}(A,A)[1]
 \]
 is a complete dg Lie algebra and the curved \(A_\infty\)-chiral
 structure is an element \(m\in\mathfrak g_A^1\) satisfying
 \[
 d_{\mathcal P}m+\frac12[m,m]=0;
 \]
\item a cofibrant diagonal \(A\)-bimodule resolution and a continuous
 cochain map
 \[
 \gamma_A\colon C^\bullet_{\mathcal P}(A,A)
 \longrightarrow R\operatorname{End}_{A^e}(A);
 \]
\item finite conformal-weight quotients, an identification of the
 complete cochain object with their inverse limit, and the derived-limit
 condition
 \[
  \varprojlim\nolimits_N^1 H^{q-1}(C_{\mathcal P,\leq N})=0
  \qquad(q\in\mathbb Z)
 \]
 whenever a cohomology calculation passes through that limit.
\end{enumerate}
\end{definition}

The package separates four objects:
\[
 A,\qquad B_X(A),\qquad A^{\mathrm i},\qquad
 Z^{\mathrm{der}}_{\mathrm{ch}}(A).
\]
Here \(B_X(A)\) is a complete pro-conilpotent coalgebra,
\(A^{\mathrm i}\) is the quadratic coalgebra of a chosen presentation,
and the derived centre is the endomorphism object
\eqref{eq:cct-derived-centre}.  Comparison maps connect these objects;
their types already distinguish them.

\subsection{The chiral endomorphism operad}
\label{subsec:chiral-endo-operad}

\begin{definition}[Chiral endomorphism-operad chart;
\ClaimStatusConditional]
\label{def:chiral-endomorphism-operad}
Assume \(H_{\mathrm{op}}(A)\).  The symbol
\(\End_A^{\mathrm{ch}}\) denotes the chosen dg operad
\(\mathcal P_A^{\mathrm{ch}}\).  On a formal coordinate disc one may
seek a Laurent chart
\[
 \End_{A,\mathrm{Laur}}^{\mathrm{ch}}(n)
 =
 \operatorname{Hom}_{\mathrm{cts}}
 \!\left(A^{\widehat\otimes n},
 A((z_1-z_n))\cdots((z_{n-1}-z_n))\right).
\]
This collection represents the operad once a common expansion
convention and continuous substitution maps have been supplied and
shown to satisfy the operadic identities.  Those requirements are
clauses~\textup{(i)}--\textup{(ii)} of
Definition~\ref{def:cct-operadic-chart-package}.
\end{definition}

\begin{remark}[Three layers of chiral spectral variables;
\ClaimStatusProvedHere]
\label{rem:endch-spectral-variable-layering}
The geometric logarithmic \(\mathcal D\)-module model, its formal
completion along a diagonal, and a Laurent-series chart are three
typed layers.  Relative coordinates \(z_i-z_n\) belong to the second
and third layers.  The Fourier variable in a vertex-algebra
\(\lambda\)-bracket belongs to the translation-equivariant operadic
description.  A coordinate/Fourier transform relates these variables
after its expansion and sesquilinearity conventions are fixed.
\end{remark}

\begin{lemma}[Associativity of partial compositions;
\ClaimStatusConditional]
\label{lem:partial-comp-assoc}
Under \(H_{\mathrm{op}}(A)\), the partial compositions of
\(\End_A^{\mathrm{ch}}\) satisfy nested and disjoint operadic
associativity.  In the Laurent chart this conclusion holds precisely
when all iterated substitutions are defined in the same completed
expansion ring.
\end{lemma}

\begin{proof}
Nested and disjoint associativity are axioms of the dg operad in
\(H_{\mathrm{op}}(A)\).  The second sentence is their chart-level
translation: equality in two different iterated Laurent rings has
mathematical content after both expressions map to a common
completion.
\end{proof}

\subsection{The chiral Hochschild cochain complex}
\label{subsec:chiral-hoch-cochain}

\begin{definition}[Chiral Hochschild cochain complex;
\ClaimStatusDefinitional]
\label{def:chiral-hochschild-cochain-brace}
The chart-derived chiral Hochschild complex is
\[
 C^\bullet_{\mathrm{ch}}(A,A)
 :=
 R\operatorname{End}_{A^e}(A).
\]
Under \(H_{\mathrm{op}}(A)\), the chosen operadic model is denoted
\(C^\bullet_{\mathcal P}(A,A)\), with
\[
 C^\bullet_{\mathcal P}(A,A)[1]=\mathfrak g_A,
 \qquad
 d_m=d_{\mathcal P}+[m,-].
\]
The comparison \(\gamma_A\) of
Definition~\ref{def:cct-operadic-chart-package} promotes this model to
the derived centre when \(\gamma_A\) is a quasi-isomorphism.
\end{definition}

\begin{definition}[Geometric--algebraic comparison package;
\ClaimStatusDefinitional]
\label{def:cct-geometric-algebraic-package}
The package \(H_{\mathrm{geom/alg}}(A)\) consists of:
\begin{enumerate}[label=\textup{(\roman*)}]
\item a geometric logarithmic cochain complex
 \(C^\bullet_{\mathrm{ch,geom}}(A,A)\) on a specified
 Fulton--MacPherson or Ran model;
\item a continuous restriction-and-expansion chain map
 \[
 \varphi_A\colon
 C^\bullet_{\mathrm{ch,geom}}(A,A)
 \longrightarrow C^\bullet_{\mathcal P}(A,A);
 \]
\item compatibility of \(\varphi_A\) with ordered operadic
 compositions;
\item a finite-window quasi-isomorphism proof;
\item inverse-limit identifications for both complexes and vanishing
 \(\varprojlim_N^1 H^{q-1}\) on both finite-window towers.
\end{enumerate}
\end{definition}

\begin{proposition}[Geometric--algebraic comparison criterion;
\ClaimStatusConditional]
\label{prop:geometric-algebraic-hochschild}
Assume \(H_{\mathrm{op}}(A)+H_{\mathrm{geom/alg}}(A)\).  Then
\[
 \varphi_A\colon
 C^\bullet_{\mathrm{ch,geom}}(A,A)
 \xrightarrow{\sim}
 C^\bullet_{\mathcal P}(A,A)
\]
is a quasi-isomorphism of complete cochain complexes and preserves the
ordered brace operations.  If \(\gamma_A\) is also a
quasi-isomorphism, both models compute
\(Z^{\mathrm{der}}_{\mathrm{ch}}(A)\).
\end{proposition}

\begin{proof}
Clause~\textup{(iv)} gives the quasi-isomorphism on every finite
window.  Clause~\textup{(v)} identifies cohomology of the inverse
limit with the inverse limit of the finite-window cohomologies.
Clause~\textup{(iii)} gives brace compatibility.  Composition with
\(\gamma_A\) gives the final statement.
\end{proof}

\begin{remark}[Algebraic vs.\ geometric content]
\label{rem:comparison-geometric-hoch}
The formal operadic brace theorem below is independent of a curve.
The construction of \(\varphi_A\), its compatibility with residues,
and its completed quasi-isomorphism are geometric theorems.  Deligne's
logarithmic de Rham comparison computes the cohomology of an open
configuration space.  Identification with a coefficient-valued
Hochschild complex and its brace operations is the additional content
of \(H_{\mathrm{geom/alg}}(A)\).
\end{remark}

\subsection{Brace operations on chiral Hochschild cochains}
\label{subsec:brace-ops-chiral}

\begin{definition}[Operadic braces; \ClaimStatusDefinitional]
\label{def:brace-operations-chiral}
For a complete dg nonsymmetric operad \(\mathcal P\), the brace
\(f\{g_1,\ldots,g_r\}\) is the signed sum over all ordered disjoint
insertions of \(g_1,\ldots,g_r\) into the input slots of \(f\).
The signs are the Koszul signs in the suspension fixed by
\(H_{\mathrm{op}}(A)\), and the sum is continuous in the chosen
completion.
\end{definition}

\begin{example}[Single brace]
\label{ex:single-brace-chiral}
The single brace is the operadic pre-Lie product
\[
 f\{g\}=\sum_i (-1)^{\epsilon_i}f\circ_i g.
\]
Each \(\epsilon_i\) is determined by moving the suspended operation
\(g\) past the inputs preceding the \(i\)-th slot.
\end{example}

\begin{definition}[Gerstenhaber bracket; \ClaimStatusDefinitional]
\label{def:gerstenhaber-bracket-chiral}
On the shifted complex
\(\mathfrak g_A=C^\bullet_{\mathcal P}(A,A)[1]\), define
\[
 [f,g]
 :=
 f\{g\}-(-1)^{|f||g|}g\{f\}.
\]
This is a degree-zero bracket on \(\mathfrak g_A\), equivalently a
degree-\((-1)\) Gerstenhaber bracket on the unshifted cochain complex.
\end{definition}

\begin{remark}[Cup product]
\label{rem:cup-product}
For a strict binary multiplication \(\mu\), the cup product is the
standard operation obtained from \(\mu\{\!f,g\!\}\), with the
suspension sign dictated by the operadic convention.  For an
\(A_\infty\)-multiplication, the higher \(m_k\) supply its coherent
homotopies.  The brace operations, cup product, and bracket therefore
belong to one operadic structure.
\end{remark}

\subsection{Braces and the homotopy Gerstenhaber structure}
\label{subsec:brace-dg-algebra}

\begin{proposition}[Pre-Lie relation for the single brace;
\ClaimStatusProvedHere]
\label{prop:pre-lie-chiral}
The single brace associated with a dg nonsymmetric operad satisfies
the graded pre-Lie identity.
\end{proposition}

\begin{proof}
Both sides enumerate the same rooted trees with two inserted vertices.
Operadic associativity identifies the evaluations and the Koszul rule
identifies their signs.  This is the pre-Jacobi identity of
Gerstenhaber--Voronov
\cite[identity~\textup{(}9\textup{)}]{GV95}.
\end{proof}

\begin{proposition}[Full brace identity;
\ClaimStatusProvedHere]
\label{prop:full-brace-chiral}
The braces associated with a dg nonsymmetric operad satisfy the full
ordered-partition brace identity.
\end{proposition}

\begin{proof}
Both sides enumerate the same planar rooted trees, grouped according
to whether each second-stage insertion lands in an outer slot or in a
previously inserted operation.  This is the operadic brace identity
of Gerstenhaber--Voronov
\cite[identity~\textup{(}6\textup{)} and Theorem~3]{GV95}.
\end{proof}

\begin{theorem}[Operadic braces and homotopy Gerstenhaber structure;
\ClaimStatusConditional]
\label{thm:brace-dg-algebra}
\textup{[Type signature: Open quadrant; complete operadic chart
presentation; Beilinson level \(3\); hypothesis package
\(H_{\mathrm{op}}(A)\).]}

Assume \(H_{\mathrm{op}}(A)\).  The underlying complete graded object
\(C^\bullet_{\mathcal P}(A,A)\) carries the operadic braces.  The
Maurer--Cartan twist \(d_m=d_{\mathcal P}+[m,-]\) makes
\(\mathfrak g_A=C^\bullet_{\mathcal P}(A,A)[1]\) a complete dg Lie
algebra.  Planar trees whose labelled vertices are cochains and whose
unlabelled vertices are components of \(m\) define the completed
minimal-operad action of Kontsevich--Soibelman
\cite[\S\S5 and~7]{KS00}.  For a strict binary multiplication this action is
the homotopy Gerstenhaber structure of Gerstenhaber--Voronov
\cite[Theorem~3 and Corollary~4]{GV95}.  Its cohomology is a
Gerstenhaber algebra.
\end{theorem}

\begin{proof}
The brace identities follow from
Propositions~\ref{prop:pre-lie-chiral}
and~\ref{prop:full-brace-chiral}.  The Maurer--Cartan equation gives,
for \(f\in\mathfrak g_A\),
\[
 d_m^2f
 =
 \left[d_{\mathcal P}m+\frac12[m,m],f\right]
 =0.
\]
Both \(d_{\mathcal P}\) and \([m,-]\) are derivations of the shifted
Lie bracket, so their sum \(d_m\) is a Lie differential.  The
planar-tree differential splits
unlabelled vertices and applies \(d_{\mathcal P}\) to their labels;
the Maurer--Cartan equation identifies the resulting sum with zero.
This is the chain-map identity for the minimal-operad action.  In the
strict binary case it specializes to identities~\textup{(6)--(8)} of
Gerstenhaber--Voronov~\cite{GV95}.  Completion is legitimate by the
continuity clause of \(H_{\mathrm{op}}(A)\).
\end{proof}

\subsubsection{Operadic proof of the brace identity}
\label{subsubsec:proof-brace-identity}
The proof of Proposition~\ref{prop:full-brace-chiral} is the rooted-tree
bijection described there.  Operadic associativity and the Koszul rule
identify the evaluations and their signs.

\subsubsection{The square of the twisted differential}
\label{subsubsec:proof-delta-squared}
The shifted multiplication \(m\) has odd dg-Lie degree \(1\).
Consequently the graded Jacobi identity gives
\[
 [m,[m,f]]=\frac12[[m,m],f],
\]
and the displayed calculation in the proof of
Theorem~\ref{thm:brace-dg-algebra} proves \(d_m^2=0\).

\begin{remark}[The Maurer--Cartan identity is the Jacobi input]
\label{rem:half-jacobi}
The square-zero statement uses the Maurer--Cartan equation in the
shifted dg Lie algebra.  This convention keeps the \(A_\infty\)
structure odd and removes the parity ambiguity produced by mixing
shifted and unshifted brace degrees.
\end{remark}

\subsubsection{Differential compatibility}
\label{subsubsec:proof-diff-compat}
The minimal-operad differential and the evaluation of unlabelled
vertices by \(m\) form the two sides of the chain-map identity.  Terms
with two copies of \(m\) combine with the internal differential into
\(d_{\mathcal P}m+\frac12[m,m]\).  For a strict multiplication,
Gerstenhaber--Voronov identity~\textup{(8)} expresses the differential
of a higher brace as the internal differential terms together with
the adjacent cup-boundary terms.

\subsection{Local chiral Swiss-cheese pairs}
\label{subsec:local-swiss-cheese}

\begin{definition}[Local chiral action datum;
\ClaimStatusDefinitional]
\label{def:local-swiss-cheese-pair}
A local chiral action datum over \(A\) consists of a complete cochain
complex \(B\) carrying an action of the same completed minimal operad
as in Theorem~\ref{thm:brace-dg-algebra}, together with continuous
evaluation maps
\[
 \mu_{1;n}\colon B\widehat\otimes A^{\widehat\otimes n}
 \longrightarrow A((z_1-z_n))\cdots((z_{n-1}-z_n))
\]
such that:
\begin{enumerate}[label=\textup{(\roman*)}]
\item the arity components assemble to a continuous graded map
 \(\Phi_\mu\colon B\to C^\bullet_{\mathcal P}(A,A)\);
\item \(\Phi_\mu d_B=d_m\Phi_\mu\);
\item for every planar-tree operation \(T\) in the completed minimal
operad,
 \[
 \Phi_\mu\bigl(T_B(b_1,\ldots,b_r)\bigr)
 =
 T_{C^\bullet_{\mathcal P}(A,A)}\bigl(
 \Phi_\mu(b_1),\ldots,\Phi_\mu(b_r)\bigr).
 \]
\end{enumerate}
The package \(H_{\mathrm{SC}}(B,A)\) is the assertion that these maps
and identities have been constructed in the same expansion and
completion conventions as \(H_{\mathrm{op}}(A)\).
\end{definition}

\begin{remark}[Relation to the operad-theoretic definition]
\label{rem:swiss-cheese-operad}
An algebra over a two-coloured Swiss-cheese chain operad supplies such
evaluation maps after restriction to one local closed disc and one
local open disc.  Establishing this restriction for a proposed chiral
Swiss-cheese operad is a separate comparison theorem.
\end{remark}

\subsection{The universal pair: construction}
\label{subsec:universal-pair-construction}

\begin{construction}[The cochain action pair;
\ClaimStatusConditional]
\label{const:universal-pair}
Under \(H_{\mathrm{op}}(A)\), operadic cochains act on \(A\) by
evaluation.  This defines the local pair
\[
 U(A):=\bigl(C^\bullet_{\mathcal P}(A,A),A,\operatorname{ev}\bigr).
\]
\end{construction}

\subsection{The chiral Deligne--Tamarkin criterion}
\label{subsec:chiral-deligne-tamarkin}

\begin{theorem}[Evaluation criterion for a local closed-sector action;
\ClaimStatusConditional]
\label{thm:chiral-deligne-tamarkin}
\textup{[Type signature: Open quadrant; local two-coloured action
presentation; Beilinson level \(3\); hypothesis package
\(H_{\mathrm{op}}(A)+H_{\mathrm{SC}}(B,A)\).]}

Assume \(H_{\mathrm{op}}(A)+H_{\mathrm{SC}}(B,A)\).  The mixed
operations define a unique continuous morphism of completed
minimal-operad algebras
\[
 \Phi_\mu\colon B\longrightarrow C^\bullet_{\mathcal P}(A,A)
\]
whose arity-\(n\) component is the specified map \(\mu_{1;n}\).
If \(\gamma_A\) is a quasi-isomorphism, this morphism represents the
induced action on \(Z^{\mathrm{der}}_{\mathrm{ch}}(A)\).
\end{theorem}

\begin{proof}
Existence is clauses~\textup{(i)}--\textup{(iii)} of
\(H_{\mathrm{SC}}(B,A)\).  A cochain in the chosen operadic model is
determined by its arity components, so the equations
\(\Phi_\mu(b)_n=\mu_{1;n}(b;-)\) force uniqueness.  Composition with
\(\gamma_A\) gives the derived-centre action.
\end{proof}

\begin{remark}[Meaning of the theorem]
The theorem is an evaluation criterion.  A terminal-object statement
in an \(\infty\)-category of chiral Swiss-cheese algebras requires a
definition of that category, coherent higher morphisms, and a proof of
the corresponding mapping-space equivalence.  Those data form an open
comparison problem.  Voronov's Swiss-cheese operad
\cite{Voronov99} and the operadic Deligne theorem of
Kontsevich--Soibelman~\cite{KS00} provide the classical templates.
\end{remark}

\begin{remark}[Koszul sign in the evaluation map;
\ClaimStatusDefinitional]
\label{rem:phi-koszul-sign}
The suspension convention in \(H_{\mathrm{op}}(A)\) determines the
sign relating \(\mu_{1;n}\) to the arity-\(n\) component of
\(\Phi_\mu\).  A universal factor such as \((-1)^{|b|}\) becomes a
theorem after the source and target suspension functors have been
fixed.
\end{remark}

\begin{definition}[E3 topologization datum for a derived chiral centre;
\ClaimStatusDefinitional]
\label{def:e3-topologization-datum-derived-centre}
An \(E_3\)-topologization datum consists of a completed BRST complex
\((\mathcal B,Q)\), a minimal-operad-compatible map from
\(Z^{\mathrm{der}}_{\mathrm{ch}}(A)\), a degree-\((-1)\) operator
\(G\), and coherent comparison data satisfying
\[
 [Q,G]=L_{-1}.
\]
The package also includes the framed-discs or Dunn-additivity
comparison that converts the null-homotopy of translation into the
extra topological direction.
\end{definition}

\begin{proposition}[E2 centre theorem versus E3 topologization;
\ClaimStatusConditional]
\label{prop:e2-center-not-e3-topologization}
Assume the datum of
Definition~\ref{def:e3-topologization-datum-derived-centre}.  For every
\(Q\)-closed observable \(x\),
\[
 L_{-1}x=Q(Gx).
\]
Thus translation acts trivially on BRST cohomology.  The supplied
framed-discs comparison then promotes the minimal-operad
\(E_2\)-type structure to the stated \(E_3\)-topological structure.
\end{proposition}

\begin{proof}
The displayed identity is the graded commutator equation
\([Q,G]=L_{-1}\) evaluated on a \(Q\)-closed element.  The final
statement is exactly the framed-discs comparison included in the
topologization datum.
\end{proof}

\subsection{Additional hypotheses and corollaries}
\label{subsec:additional-hypotheses}

\begin{remark}[Faithfulness and injectivity]
\label{rem:faithfulness-hypothesis}
For a local action datum, \(\ker\Phi_\mu\) consists of the closed
operations whose evaluations on all open inputs vanish.  The map is
injective precisely when the open action is faithful.
\end{remark}

\begin{corollary}[The chiral center;
\ClaimStatusConditional]
\label{cor:chiral-center}
Assume \(H_{\mathrm{op}}(A)\) and that \(\gamma_A\) is a
quasi-isomorphism.  Then
\[
 Z_{\mathrm{ch}}(A)
 :=
 H^0\!\left(Z^{\mathrm{der}}_{\mathrm{ch}}(A)\right)
 \cong H^0\!\left(C^\bullet_{\mathcal P}(A,A)\right).
\]
The minimal-operad homotopy for the cup product makes this degree-zero
cohomology a commutative algebra.
\end{corollary}

\begin{proof}
The quasi-isomorphism gives the displayed identification.  The
operadic brace homotopy identifies the two orders of the cup product
on cohomology.
\end{proof}

\begin{corollary}[Gerstenhaber structure on cohomology;
\ClaimStatusConditional]
\label{cor:gerstenhaber-cohomology}
Under the same hypotheses,
\(H^\bullet(Z^{\mathrm{der}}_{\mathrm{ch}}(A))\) carries the
Gerstenhaber structure transported from the operadic cochain model.
\end{corollary}

\begin{remark}[Comparison with the classical Deligne conjecture]
\label{rem:classical-deligne}
For an associative algebra, Hochschild cochains carry an \(E_2\)
structure.  Kontsevich--Soibelman~\cite{KS00} give an operadic model,
and Voronov~\cite{Voronov99} supplies the Swiss-cheese geometry.  The
chiral theorem requires the additional construction of
\(H_{\mathrm{op}}(A)\), the geometric comparison
\(H_{\mathrm{geom/alg}}(A)\), and any desired two-coloured comparison.
\end{remark}

\begin{remark}[Bridge to the global open/closed theory]
\label{rem:chiral-center-global-bridge}
A global open/closed theory produces a local action on the chart
derived centre once restriction to a formal disc, compatibility with
the chart resolution, and completion have been supplied.  These are
the local-global and OCA comparison packages.  The resulting map is
the morphism of
Theorem~\ref{thm:chiral-deligne-tamarkin}; a physical bulk
interpretation belongs to the additional open/closed realization.
\end{remark}


%% ===========================================================
%% EXPLICIT DERIVED CENTER COMPUTATIONS
%% ===========================================================

\subsection{Bounded computations and derived-centre promotion}
\label{subsec:derived-center-explicit}
\index{derived center!bounded comparison|textbf}

The chapter uses two cochain functors.  Bakalov--De Sole--Kac form the
bounded vertex complex
\(C^\bullet_{\mathrm{ch},b}(V,V)\) from the bounded part of the
translation-equivariant chiral operad.  The chart-derived centre is
\[
 C^\bullet_{\mathrm{ch}}(\cA_V,\cA_V)
 =
 R\operatorname{End}_{\cA_V^e}(\cA_V).
\]
A statement about the first complex becomes a statement about the
second through a specified bounded-to-chart quasi-isomorphism.

\begin{definition}[Finite-window promotion package;
\ClaimStatusDefinitional]
\label{def:cct-finite-window-promotion}
For a vertex algebra \(V\) and its curve chart \(\cA_V\), the package
\(H_{\mathrm{fw}}(V)\) consists of:
\begin{enumerate}[label=\textup{(\roman*)}]
\item a cochain map
 \[
 \chi_V^{\mathrm{bd}}\colon
 C^\bullet_{\mathrm{ch},b}(V,V)
 \longrightarrow
 C^\bullet_{\mathrm{ch}}(\cA_V,\cA_V);
 \]
\item a proof that \(\chi_V^{\mathrm{bd}}\) is a quasi-isomorphism on
 each finite conformal-weight window;
\item surjective transition maps and vanishing \(\varprojlim^1\) for
 the completed tower;
\item compatibility with any product, bracket, cyclic operator, or
 family parameter used in the promoted conclusion.
\end{enumerate}
\end{definition}

\begin{proposition}[Bounded benchmarks and chart criterion;
\ClaimStatusConditional]
\label{prop:derived-center-explicit}
\begin{enumerate}[label=\textup{(\roman*)}]
\item For a one-dimensional even space \(\mathfrak h\),
Bakalov--De Sole--Kac compute
\[
 \dim H^n_{\mathrm{ch},b}(B_{\mathfrak h},B_{\mathfrak h})
 =(2,1,0,0,\ldots)
 \qquad\text{\cite[Theorem~7.4]{BDSK21}}.
\]
\item For the universal Virasoro vertex algebra, for every central
charge \(c\), they compute
\[
 \dim H^n_{\mathrm{ch},b}(\operatorname{Vir}_c,\operatorname{Vir}_c)
 =
 \begin{cases}
 1,&n\in\{0,2,3\},\\
 0,&n\in\mathbb Z_{\ge0}\setminus\{0,2,3\},
 \end{cases}
 \qquad\text{\cite[Theorem~7.2]{BDSK21}}.
\]
\item Under \(H_{\mathrm{fw}}(V)\), either bounded calculation
transports to
\(\ChirHoch^\bullet(\cA_V)\).
\item Universal affine vertex-algebra cohomology retains the status of
\cite[Conjecture~7.5]{BDSK21}.  Principal
\(\mathcal W\)-algebras require their own cochain model and
bounded-to-chart comparison.
\end{enumerate}
\end{proposition}

\begin{proof}
The first two clauses are the cited computations.  In the first,
\(\Pi\mathfrak h\) is an odd line, so its super-symmetric powers occur
in degrees zero and one.  The third clause is invariance of cohomology
under the quasi-isomorphism and inverse-limit comparison in
\(H_{\mathrm{fw}}(V)\).  The fourth records the status in the cited
source and isolates the comparison theorem required for the
additional families.
\end{proof}

\begin{proposition}[Affine inner-direction firewall;
\ClaimStatusProvedHere]
\label{prop:chirhoch1-affine-km}
Let \(V_k(\fg)\) be an affine vertex algebra.  For \(X\in\fg\), the
zero mode
\[
 \delta_X=(J^X)_{(0)}
\]
is an inner chiral derivation, and
\[
 [\delta_X,\delta_Y]=\delta_{[X,Y]}.
\]
Consequently the adjoint current directions have zero image in
\[
 \operatorname{Der}_{\mathrm{ch}}(V_k(\fg))/
 \operatorname{Inn}_{\mathrm{ch}}(V_k(\fg)).
\]
On a surface carrying an exhaustion theorem for all continuous outer
derivations, this quotient equals zero.
\end{proposition}

\begin{proof}
The simple-pole coefficient of the affine OPE gives
\((J^X)_{(0)}J^Y=J^{[X,Y]}\).  Zero modes therefore act by inner
derivations and satisfy the displayed commutator identity.  Passing to
the quotient kills their classes.
\end{proof}

\begin{remark}[Affine zero-mode firewall]
\label{rem:sl2-chirhoch-dim5}
The adjoint vector space in the current algebra and in the first
bar-dual generator is a prequotient.  The remaining continuous
derivations and the bounded-to-chart comparison determine the first
chart Hochschild group.
\end{remark}

\begin{proposition}[Generic affine second-cohomology criterion;
\ClaimStatusConditional]
\label{prop:chirhoch2-affine-km-general}
Let \(V_k(\fg)\) lie on a generic finite-window PBW surface.  Assume a
cofibrant diagonal resolution, a contraction of every positive-weight
complement, strict Mittag--Leffler completion, and a chain
identification of the degree-two endpoint with the vacuum line of the
specified dual chart.  Then
\[
 \ChirHoch^2(V_k(\fg))\cong\mathbb C.
\]
The Kac--Shapovalov determinant and the scalar centre of a dual affine
algebra provide inputs to this package.  The diagonal chain map,
contraction, and completion supply the remaining data needed to
compute the diagonal Ext group.
\end{proposition}

\begin{proof}
The assumed contraction removes the positive-weight complement.  The
endpoint comparison identifies the remaining degree-two cohomology
with the one-dimensional vacuum line, and the Mittag--Leffler clause
passes the calculation to the completion.
\end{proof}

\begin{remark}[Finite-window Kac--Shapovalov test]
\label{rem:chirhoch2-affine-km-general-falsification}
A finite-window determinant calculation tests the vacuum-centre input.
A proof of Proposition~\ref{prop:chirhoch2-affine-km-general} also
needs the diagonal chain map, its contraction, and completion.
\end{remark}

\begin{proposition}[Affine Gerstenhaber bracket on inner directions;
\ClaimStatusProvedHere]
\label{prop:gerstenhaber-sl2-bracket}
The prequotient zero-mode derivations satisfy
\[
 [\delta_X,\delta_Y]=\delta_{[X,Y]}.
\]
Their images in first Hochschild cohomology are zero, so the induced
Gerstenhaber bracket on this inner subspace is zero.  Additional outer
classes form a separate cohomological sector.
\end{proposition}

\begin{proof}
This is Proposition~\ref{prop:chirhoch1-affine-km} together with
functoriality of the bracket under passage to cohomology.
\end{proof}

\begin{remark}[First non-abelian zero-mode bracket verification]
\label{rem:first-nonabelian-gerstenhaber}
For \(\fg=\mathfrak{sl}_2\), the current bracket is the first
non-abelian prequotient example.  Its image vanishes because every
current zero mode is inner.
\end{remark}

\begin{definition}[DS--Hochschild comparison package;
\ClaimStatusDefinitional]
\label{def:cct-ds-hochschild-package}
The package \(H_{\mathrm{DS/HH}}(k,\fg)\) consists of a functorial
BRST model for Drinfeld--Sokolov reduction, a chain map between
cofibrant diagonal resolutions, compatibility with operadic braces,
and convergent finite-window completion.
\end{definition}

\begin{proposition}[DS--ChirHoch compatibility criterion;
\ClaimStatusConditional]
\label{prop:ds-chirhoch-compatibility}
Under \(H_{\mathrm{DS/HH}}(k,\fg)\), Drinfeld--Sokolov reduction
induces a map on chart chiral Hochschild cohomology.  The affine
current zero-mode directions map to zero because their source classes
already vanish in the inner quotient.
\end{proposition}

\begin{proof}
The chain map in \(H_{\mathrm{DS/HH}}\) induces the cohomology map.
Proposition~\ref{prop:chirhoch1-affine-km} identifies each current
zero-mode class with zero before applying that map.
\end{proof}

\begin{remark}[Generalization to \(\mathfrak{sl}_N\)]
\label{rem:ds-chirhoch-sl-n}
The same inner-direction argument applies to every simple \(\fg\).
A statement about all first cohomology classes of the DS quotient
requires the full package \(H_{\mathrm{DS/HH}}(k,\fg)\) and a
calculation of the target complex.
\end{remark}

\begin{conjecture}[Heisenberg cyclic/BV enhancement;
\ClaimStatusConjectured]
\label{prop:heisenberg-bv-structure}
A cyclic comparison for the Heisenberg chart is expected to transport
the Connes operator to a BV operator on its chart-derived centre.
A proof requires:
\begin{enumerate}[label=\textup{(\roman*)}]
\item a bounded or finite-window cyclic cochain model;
\item a quasi-isomorphism to the chart cyclic complex;
\item compatibility of the Connes operator with completion;
\item a direct calculation of products and brackets on representatives.
\end{enumerate}
The bounded dimension vector of
Proposition~\ref{prop:derived-center-explicit}(i) supplies the graded
vector-space input.  The four comparison clauses above determine the
cyclic structure constants.
\end{conjecture}

\begin{conjecture}[Chiral HKR comparison problem;
\ClaimStatusConjectured]
\label{prop:chirhoch-cdr}
Let \(X\) be a smooth variety and
\(\Omega_X^{\mathrm{ch}}\) its chiral de Rham sheaf.  A chiral HKR
theorem would consist of a natural quasi-isomorphism from the selected
chart Hochschild complex of a commutative chiral structure sheaf to a
complex built from \(\Omega_X^{\mathrm{ch}}\), together with its
descent and completion data.
\end{conjecture}

\begin{remark}[Attribution]
\label{rem:chirhoch-cdr-attribution}
Malikov--Schechtman--Vaintrob construct
\(\Omega_X^{\mathrm{ch}}\) as a sheaf of vertex superalgebras
\cite{MSV99}.  That construction supplies the proposed target.  The
chart Hochschild quasi-isomorphism of
Conjecture~\ref{prop:chirhoch-cdr} is the additional comparison
problem.
\end{remark}

\begin{proposition}[Annulus-trace criterion;
\ClaimStatusConditional]
\label{prop:annulus-trace-standard}
Assume a cyclic chart model, a compatible annulus trace
\[
 \operatorname{Tr}_{\cA}\colon
 \ChirHoch_*(\cA)\longrightarrow
 H^*(\overline{\mathcal M}_{1,1}),
\]
and the family normalization \(H_{\mathrm{mod}}(\cA)\).  Then its
vacuum value has the form
\[
 \operatorname{Tr}_{\cA}(1)=\kappa(\cA)\lambda_1.
\]
The coefficient is determined on families carrying both a genus-one
curvature calculation and the trace comparison.
\end{proposition}

\begin{proof}
This is the defining scalar projection in the stated cyclic and
modular comparison package.  Family values enter through their
independently verified genus-one curvature calculations.
\end{proof}

\subsection{Normalized scalar traces of Theorem~C eigensummands}
\label{subsec:theorem-C-strengthened}

The scalar complementarity problem begins with the C1 eigensummands of
$\mathbf C_g(A)=R\Gamma(\overline{\mathcal M}_g,\mathcal Z(A))$.
It asks for a genus-one characteristic \(\kappa(A)\), a specified
Verdier--Koszul partner \(A^!\), and a Theorem~D trace comparison that
transports both characteristics into the same scalar normalization.

\begin{definition}[Scalar complementarity package;
\ClaimStatusDefinitional]
\label{def:cct-scalar-complementarity-package}
The package \(H_C(A,A^!)\) consists of:
\begin{enumerate}[label=\textup{(\roman*)}]
\item the C0 strict-flat fibre--centre comparison defining the ordinary
 centre local system~$\mathcal Z(A)$;
\item the represented C1 involution~$\sigma$ and its projectors
 $p^\pm$ on
 $\mathbf C_g(A)=R\Gamma(\overline{\mathcal M}_g,\mathcal Z(A))$;
\item the continuous Verdier comparison defining \(A^!\);
\item genus-one curvature computations defining
 \(\kappa(A)\) and \(\kappa(A^!)\);
\item diagonal degree-two representatives in the two C1 eigensummands
 and the Theorem~D trace maps whose normalized values are
 $\kappa(A)$ and~$\kappa(A^!)$;
\item central-charge normalizations \(c(A)\), \(c(A^!)\), together with
 a theorem identifying the family involution and the two trace
 conventions;
\item completion and averaging maps carrying these data to one scalar
 lane.
\end{enumerate}
When \(c(A)\) is invertible, write
\[
 \varrho(A):=\frac{\kappa(A)}{c(A)},\qquad
 K^c(A,A^!):=c(A)+c(A^!),\qquad
 K^\kappa(A,A^!):=\kappa(A)+\kappa(A^!).
\]
\end{definition}

\begin{theorem}[Computed and open scalar complementarity lanes;
\ClaimStatusConditional]
\label{thm:scalar-trace-complementarity}
\label{thm:derived-centre-complementarity-strengthened}
\textup{[Type signature: Open quadrant; symmetric scalar trace
presentation; Beilinson level \(5\); hypothesis package
\(H_C(A,A^!)\).]}

Under~$H_C(A,A^!)$, the quantity~$K^\kappa(A,A^!)$ is the sum of the
two normalized Theorem~D trace outputs on the C1 eigensummands.

\textup{(a)} If \(H_C(A,A^!)\) identifies
\(\varrho(A)=\varrho(A^!)=\varrho\), then
\begin{equation}\label{eq:thm-C-bridge}
 K^\kappa(A,A^!)=\varrho\,K^c(A,A^!).
\end{equation}

\textup{(b)} The unconditional Vol~I landmark values are
\[
 \mathcal S^{\mathrm{VolI,exact}}_\kappa=\{0,\;13\}.
\]
The package
$H_{\mathrm{diag}}^{g=1}+H_{W_3}^{\mathrm{DS/bar}}$ supplies the
conditional principal-$\mathcal W_3$ extension
\begin{equation}\label{eq:thm-C-ceiling}
 \mathcal S^{\mathrm{VolI,comparison}}_\kappa
 =
 \left\{0,\;13,\;\frac{250}{3}\right\}.
\end{equation}
The status-separated calculations are:
\[
\begin{array}{c|c|c|c}
\text{lane} & K^c & K^\kappa & \text{status}\\
\hline
\text{free/affine cancellation} & \text{family-specific} & 0 &
\text{computed}\\
\operatorname{Vir}_c & 26 & 13 & \text{computed}\\
\mathcal W_3^{\mathrm{prin}} & 100 & 250/3 &
H_{\mathrm{diag}}^{g=1}+H_{W_3}^{\mathrm{DS/bar}}
\end{array}
\]
The corresponding anomaly ratios are family-specific on the first row,
$1/2$ on the Virasoro row, and conditionally $5/6$ on the principal
$\mathcal W_3$ row.

\textup{(c)} For the standard Fehily--Kawasetsu--Ridout conformal
vector of the Bershadsky--Polyakov algebra
\cite[Definition~2.1]{FKR20}, for \(k\in\mathbb C\setminus\{-3\}\),
\[
 c_{\mathrm{BP}}(k)
 =
 -\frac{(2k+3)(3k+1)}{k+3},
 \qquad
 c_{\mathrm{BP}}(k)+c_{\mathrm{BP}}(-k-6)=50.
\]
Thus \(K^c_{\mathrm{BP}}=50\) is exact.  The strong generators
\(J,G^+,G^-,T\) are even, and their reciprocal-weight diagnostic is
\[
 1+\frac23+\frac23+\frac12=\frac{17}{6}.
\]
The characteristic \(\kappa_{\mathrm{BP}}\), the anomaly ratio
\(\varrho_{\mathrm{BP}}\), and
\(K^\kappa_{\mathrm{BP}}\) have status \ClaimStatusOpen{}: their
genus-one computation is the open obligation.  The equation
\(\kappa_{\mathrm{BP}}=c_{\mathrm{BP}}/6\), hence
\(K^\kappa_{\mathrm{BP}}=25/3\), is retained as the explicit
consequence of that conditional genus-one equation.

\textup{(d)} The Mukai lattice of a K3 surface is
\[
 \widetilde H(K3,\mathbb Z)
 \cong U^4\oplus E_8(-1)^2
 =\mathrm{II}_{4,20}.
\]
It has rank \(24\), signature \((4,20)\), and the exact arithmetic
value
\[
 2c_+(\mathrm{II}_{4,20})=8.
\]
This is the Mukai-lattice calculation of \cite{Mukai1987}.
Its interpretation as
\(K^\kappa_{\mathsf B}=8\) is conjectural under
\[
 H_B=
 H_{\mathrm{chart}}+H_{\mathrm{KD}}+H_{\mathrm{scalar}}
 +H_{\mathrm{mod}}+H_{\mathrm{quantum}}.
\]
The number \(8\) is therefore a lattice theorem and a chiral
complementarity candidate.

\textup{(e)} For a principal \(\mathcal W_N\) family, the formula
\begin{equation}\label{eq:thm-C-WN}
 K^\kappa(\mathcal W_N,\mathcal W_N^!)
 =(H_N-1)K^c(\mathcal W_N,\mathcal W_N^!)
\end{equation}
is a conditional family comparison for \(N\ge3\): it requires the
genus-one characteristic, Verdier partner, and common ratio theorem
in \(H_C\).  The \(N=3\) entry in
\eqref{eq:thm-C-ceiling} is its first conditional landmark.
\end{theorem}

\begin{proof}
Part~\textup{(a)} is the calculation
\[
 \kappa(A)+\kappa(A^!)
 =\varrho c(A)+\varrho c(A^!)
 =\varrho K^c(A,A^!).
\]
For Virasoro,
\(c/2+(26-c)/2=13\).  For principal \(\mathcal W_3\),
the exact central-charge reflection gives~\(100\); on
$H_{\mathrm{diag}}^{g=1}+H_{W_3}^{\mathrm{DS/bar}}$, the ratio
\(5/6\) gives \(250/3\).  The free and affine lanes pair opposite scalar
characteristics and give zero.

For BP, direct substitution under \(k\mapsto-k-6\) proves the
central-charge identity.  The parity statement gives the displayed
positive reciprocal-weight sum.  A genus-one curvature calculation
with charged ghosts, neutral fields, improvement term, and mixed
channels is the proof obligation for the three open invariants.

For K3, the block decomposition
\(U^4\oplus E_8(-1)^2\), the Betti-number/signature calculation, and
the explicit even unimodular Gram matrix give independent proofs of
rank \(24\) and signature \((4,20)\).  The passage from
\(2c_+=8\) to a chiral conductor is exactly the package \(H_B\).
Part~\textup{(e)} is the same bridge calculation under its stated
family hypotheses.
\end{proof}

\begin{equation}\label{eq:thm-C-full-set}
 \mathcal S^\mathrm{comparison}_\kappa
 =
 \mathcal S^{\mathrm{VolI,exact}}_\kappa
 \;\sqcup\;
 \left\{\frac{250}{3}\right\}_{H_{\mathrm{diag}}^{g=1}+H_{W_3}^{\mathrm{DS/bar}}}
 \;\sqcup\;
 \left\{\frac{25}{3}\right\}_{\kappa_{\mathrm{BP}}=c_{\mathrm{BP}}/6}
 \;\sqcup\;
 \{8\}_{H_B\text{-candidate}}.
\end{equation}
The status tags in this display are part of the notation: the first set
comprises the unconditional Vol~I lanes; each further entry records the
hypothesis package that produces it.  The Mukai entry is a conjectural
chiral interpretation of exact lattice arithmetic.

\begin{remark}[Test families and the scalar frontier]
\label{rem:theorem-C-test-families}
Lattice and free-field families test the zero lane.  Virasoro tests
the value \(13\), and principal \(\mathcal W_3\) tests the conditional
value \(250/3\) under its diagonal-trace and DS/bar packages.
Bershadsky--Polyakov tests the distinction between an exact
central-charge conductor and an open modular characteristic.  The K3
Mukai lattice tests the distinction between exact lattice arithmetic
and a conjectural chiral comparison.  Principal \(\mathcal W_N\) for
\(N\ge4\) tests the additional family package in
\eqref{eq:thm-C-WN}.
\end{remark}

\section{Open automorphic realization problem}
\label{sec:complementarity-baily-borel}

\begin{remark}[Baily--Borel/Humbert comparison datum;
\ClaimStatusConjectured]
\label{rem:cct-baily-borel}
The Baily--Borel stratification
\[
 \overline{\mathcal A}_2^{\mathrm{BB}}
 =
 \mathcal A_2\sqcup\mathcal A_1\sqcup\mathcal A_0,
 \qquad \mathcal A_0=\{\mathrm{cusp}\},
\]
is the genus-two Siegel boundary stratification of
Baily--Borel~\cite{BailyBorel66}.  A realization of the scalar
complementarity lanes on these strata would require:
\begin{enumerate}[label=\textup{(\roman*)}]
\item a functor from each specified chiral family to an automorphic
 or Hodge-theoretic object on \(\overline{\mathcal A}_2^{\mathrm{BB}}\);
\item a theorem identifying the chiral genus-one trace with a divisor,
 Chern class, or boundary specialization on the target;
\item compatibility with a specified Humbert-divisor inclusion and cusp
 specialization;
\item normalization checks distinguishing \(K^c\), \(K^\kappa\),
 automorphic weight, and the Mukai lattice number.
\end{enumerate}
The present status is an open comparison problem.  The scalar values
\(0\), \(13\), and \(250/3\) arise from their stated chiral
calculations and comparison packages; BP and Mukai retain the statuses recorded in
Theorem~\ref{thm:scalar-trace-complementarity}.
\end{remark}
